\section*{Chapitre \ref{ch:g12:contdist} --- Variables aléatoires continues}

\begin{solution}{exo:g12:contdist:1}
$\P(X \leq 5) = \frac{5}{15} = \frac13$~;
$\P(4 \leq X \leq 10) = \frac{6}{15} = \frac25$~;
$\E(X) = \frac{15}{2} = 7.5$~min~;
$\sigma(X) = \frac{15}{\sqrt{12}} = \frac{5\sqrt3}{2} \approx 4.3$~min.
\end{solution}

\begin{solution}{exo:g12:contdist:2}
$f \geq 0$ et
$\int_0^2 \frac38 t^2\,\dd t = \frac38\left[\frac{t^3}{3}\right]_0^2
= \frac38 \times \frac83 = 1$~: une \omterm{def:g12:contdist:density}{densité}.
\[
\E(X) = \int_0^2 \frac38 t^3\,\dd t = \frac38 \times \frac{16}{4} = \frac32,
\qquad
\P(X \geq 1) = \int_1^2 \frac38 t^2\,\dd t = \frac{8 - 1}{8} = \frac78 .
\]
\end{solution}

\begin{solution}{exo:g12:contdist:3}
\emph{1.} $\E(X) = \frac{1}{0.2} = 5$~ans~;
$\P(X > 5) = \eu^{-0.2\times5} = \eu^{-1} \approx 0.37$.

\emph{2.} Par l'\omterm{thm:g12:contdist:memoryless}{absence de mémoire} (\cref{thm:g12:contdist:memoryless}),
$\pcond{X > 3}{X > 8} = \P(X > 5) = \eu^{-1} \approx 0.37$~: les trois ans
de service ne changent rien.
\end{solution}

\begin{solution}{exo:g12:contdist:4}
$\P(X > m) = \eu^{-\lambda m} = \frac12$ donne
$m = \frac{\ln 2}{\lambda} \approx \frac{0.69}{\lambda}$, plus petite que
$\E(X) = \frac1\lambda$. La \omterm{def:g12:contdist:density}{densité} exponentielle a une longue queue à
droite~: quelques durées de vie atypiquement longues tirent la
\emph{\omterm{def:g11:stat:mean}{moyenne}} au-dessus de la \emph{\omterm{def:g11:stat:median}{médiane}}, la valeur dépassée par la
moitié de la population. (C'est la demi-vie de
l'\cref{ex:g12:exp:halflife}~: la moitié des atomes y survivent, bien que
la durée de vie \omterm{def:g11:stat:mean}{moyenne} soit plus longue.)
\end{solution}

\begin{solution}{exo:g12:contdist:5}
$168 = \mu - \sigma$ et $182 = \mu + \sigma$~: environ $68\%$.
$189 = \mu + 2\sigma$~: au-dessus se trouve la moitié du reste
$100 - 95.4 = 4.6\%$, donc environ $2.3\%$.
$154 = \mu - 3\sigma$~: environ $\frac{100 - 99.7}{2} = 0.15\%$.
\end{solution}

\begin{solution}{exo:g12:contdist:6}
$\P(X < 500) \leq 2.5\%$ signifie que $500$ doit se situer au moins deux
\omterm{def:g11:stat:variance}{écarts types} sous la \omterm{def:g11:stat:mean}{moyenne} (la normale laisse $2.3\% \approx 2.5\%$ sous
$\mu - 2\sigma$~; avec le conventionnel $1.96$, le raisonnement est
identique)~:
\[
\mu - 2\sigma \geq 500 \iff \mu \geq 500 + 8 = 508 \text{ g}.
\]
La machine doit être réglée sur $508$~g en \omterm{def:g11:stat:mean}{moyenne} --- le prix de la
garantie est $8$~g de produit «~gratuit~» par sac.
\end{solution}

\begin{solution}{exo:g12:contdist:7}
\emph{1.} Avec $p = \frac16$, $n = 1200$~:
$\sqrt{\frac{p(1-p)}{n}} = \sqrt{\frac{5/36}{1200}} \approx 0.0108$, donc
l'\omterm{def:g10:numbers:interval}{intervalle} de fluctuation à $95\%$ est
\[
\frac16 \pm 1.96 \times 0.0108 \approx \intcc{0.146}{0.188}.
\]

\emph{2.} La \omterm{def:g10:stats:series}{fréquence} observée $\frac{260}{1200} \approx 0.217$ est
largement hors~: l'hypothèse d'équilibre est rejetée au niveau $5\%$.
L'écart est d'environ $4.6$ \omterm{def:g11:stat:variance}{écarts types} --- pour une \omterm{def:g10:numbers:approx}{approximation}
normale, une \omterm{def:g10:proba:distribution}{probabilité} de l'ordre de $10^{-6}$, bien plus concluante que
la borne $\leq 4.6\%$ de Bienaymé--Tchebychev
(\cref{exo:g12:sums:7}).
\end{solution}

\begin{solution}{exo:g12:contdist:8}
\emph{1.} Une marge $\frac{1}{\sqrt n} \leq 0.02$ exige
$n \geq 2500$ personnes.

\emph{2.} Pour distinguer des scores séparés d'un point, la marge doit
être bien sous $0.5$ point, exigeant
$n \geq \left(\frac{1}{0.005}\right)^2 = 40\,000$ par la formule
simplifiée --- déjà impraticable pour la plupart des sondages. Pire, la
marge statistique ne compte que l'erreur d'\emph{échantillonnage}~; les
biais systématiques (\omterm{def:g10:proba:sample}{échantillons} non représentatifs, non-réponse, bascules
de dernière minute) ne diminuent pas avec $n$ et dépassent typiquement un
point. Une course à $1$ point est réellement trop serrée pour être
appelée.
\end{solution}

\begin{solution}{exo:g12:contdist:9}
\emph{1.} Comme $\alpha > 0$, $c \leq \alpha X + \beta \leq d
\iff \frac{c - \beta}{\alpha} \leq X \leq \frac{d - \beta}{\alpha}$, donc
\[
\P(c \leq Y \leq d)
= \int_{(c-\beta)/\alpha}^{(d-\beta)/\alpha} f(t)\,\dd t .
\]
La substitution $y = \alpha t + \beta$ (c.-à-d.\ lire l'aire dans la
variable $y$, $t = \frac{y - \beta}{\alpha}$,
$\dd t = \frac{\dd y}{\alpha}$) transforme ceci en
$\int_c^d \frac{1}{\alpha} f\left(\frac{y-\beta}{\alpha}\right)\dd y$~: la
\omterm{def:g11:func:function}{fonction} $g(y) = \frac1\alpha f\left(\frac{y-\beta}{\alpha}\right)$ est
une \omterm{def:g12:contdist:density}{densité} de $Y$.

\emph{2.} Avec $f = \varphi$, $\alpha = \sigma$, $\beta = \mu$~:
\[
g(y) = \frac{1}{\sigma}\,\varphi\!\left(\frac{y - \mu}{\sigma}\right)
= \frac{1}{\sigma\sqrt{2\pi}}\,
\exp\left(-\frac{(y-\mu)^2}{2\sigma^2}\right),
\]
qui est donc la \omterm{def:g12:contdist:density}{densité} de
$\mathcal N(\mu, \sigma^2) = \mu + \sigma\,\mathcal N(0,1)$.
\end{solution}

\begin{solution}{pb:g12:contdist:1}
\textbf{1.} $\P(\text{attente} \leq 10) = \frac{10}{30} =
\frac13$~; $\E = 15$ min~;
$\sigma = \frac{30}{\sqrt{12}} \approx 8.7$ min.

\textbf{2.} $\int_0^1 2x\,\dd x = 1$~: c'est bien une \omterm{def:g12:contdist:density}{densité}.
$\E(X) = \int_0^1 2x^2\,\dd x = \frac23$~;
$\P\left(X > \frac12\right) = \int_{1/2}^1 2x\,\dd x =
1 - \frac14 = \frac34$.

\textbf{3.} $\P(T > 15) = \eu^{-15/10} \approx 0.22$. \omterm{def:g11:stat:median}{Médiane}~:
$\eu^{-m/10} = \frac12$, donc $m = 10\ln 2 \approx 6.9$ min ---
inférieure à la \omterm{def:g11:stat:mean}{moyenne} $10$ parce que la longue traîne droite de
l'exponentielle (de rares attentes énormes) tire la \omterm{def:g11:stat:mean}{moyenne} vers le
haut tandis que la \omterm{def:g11:stat:median}{médiane} reste à l'attente «~typique~».

\textbf{4.} Par \omterm{thm:g12:contdist:memoryless}{absence de mémoire},
$\P(T > 35 \mid T > 20) = \P(T > 15) \approx 0.22$~: les vingt
minutes déjà passées n'achètent rien --- le flux ne vieillit pas.

\textbf{5.} (a) exponentielle (la désintégration radioactive est le
modèle pour lequel l'\omterm{thm:g12:contdist:memoryless}{absence de mémoire} a été inventée)~; (b)
uniforme sur $\intcc{0}{8}$ (un horaire régulier avec un décalage
aléatoire)~; (c) exponentielle~; (d) aucun des deux~: une ampoule
usée a \emph{plus} de chances de griller dans l'heure qui vient
qu'une ampoule neuve, donc
$\P(T > s + t \mid T > s) < \P(T > t)$~: le vieillissement
contredit le théorème d'\omterm{thm:g12:contdist:memoryless}{absence de mémoire}.

\textbf{6.} $68\,\%$, $95\,\%$, $99.7\,\%$~: les trois \omterm{def:g10:numbers:interval}{intervalles}
\omterm{def:g10:coordgeom:system}{repères}.

\textbf{7.} $z = \frac{190 - 172}{8} = 2.25$~: queue supérieure
d'environ $1.2\,\%$.

\textbf{8.} $z = 2$~: environ $2.3\,\%$ --- à peu près une personne
sur $44$.

\textbf{9.} La spécification se situe à $\pm 2\sigma$~: environ
$95.4\,\%$ passent et $4.6\,\%$ sont rejetés --- ruineux à grande
échelle. À $\pm 6\sigma$, le taux de défaut tombe à environ deux
pièces par \emph{milliard}~: le six sigma achète le droit de
produire en masse sans inspecter en masse.

\textbf{10.} $\sigma = 5$~; avec la correction de \omterm{def:g12:limcont:continuity}{continuité},
$\P \approx \P\left(\abs Z \leq \frac{5.5}{5}\right) =
\P(\abs Z \leq 1.1) \approx 0.73$. Le théorème lisse la planche de
Galton du \cref{pb:g11:binom:1}~: l'escalier des cases devient la
cloche \omterm{def:g12:limcont:continuity}{continue}.

\textbf{11.} $n \approx \frac{1.96^2 \times 0.25}{0.005^2}
\approx 38\,400$ personnes --- et à cette taille les erreurs
\emph{autres} que d'échantillonnage (plan de sondage, refus,
mensonges) écrasent la marge~: une course à un point d'écart est
hors de portée d'un sondage honnête, ce pour quoi les instituts
sérieux disent alors «~trop serré pour se prononcer~».

\textbf{12.} \omterm{thm:g12:contdist:memoryless}{Absence de mémoire}~: à partir de l'instant où vous
arrivez, l'attente restante est exponentielle de \omterm{def:g11:stat:mean}{moyenne} $10$ ---
les $10$ minutes pleines, et non $5$. L'«~\omterm{def:g10:numbers:interval}{intervalle} moyen de $10$~»
de l'horaire est trompeur~: votre attente suit la même \omterm{def:g12:randvar:rv}{loi} qu'un
\omterm{def:g10:numbers:interval}{intervalle} entier.

\textbf{13.} Les arrivées aléatoires tombent proportionnellement à
la longueur des \omterm{def:g10:numbers:interval}{intervalles}~: la \omterm{def:g10:proba:distribution}{probabilité} de tomber dans un
grand \omterm{def:g10:numbers:interval}{intervalle} vaut $\frac{15}{20} = \frac34$. Attente
espérée~: $\frac34 \times \frac{15}{2} + \frac14 \times
\frac52 = 5.625 + 0.625 = 6.25$ min --- plus que les $5$ naïfs,
bien que l'\omterm{def:g10:numbers:interval}{intervalle} moyen soit $10$. Coupable~:
l'\emph{échantillonnage biaisé par la longueur} --- les grands
\omterm{def:g10:numbers:interval}{intervalles} attrapent plus de voyageurs.

\textbf{14.} Les attentes vers l'arrière et vers l'avant sont
chacune exponentielles de \omterm{def:g11:stat:mean}{moyenne} $10$ (l'\omterm{thm:g12:contdist:memoryless}{absence de mémoire} joue
dans les deux sens à partir de votre arrivée)~: l'\omterm{def:g10:numbers:interval}{intervalle} qui
vous contient mesure en \omterm{def:g11:stat:mean}{moyenne} $20$ minutes --- le double de
l'\omterm{def:g10:numbers:interval}{intervalle} typique. Le paradoxe de l'inspection en une phrase~:
\emph{l'\omterm{def:g10:numbers:interval}{intervalle} que vous inspectez n'est pas un \omterm{def:g10:numbers:interval}{intervalle}
typique, parce que vous aviez plus de chances de tomber dans un
grand.}

\textbf{15.} Taille \omterm{def:g11:stat:mean}{moyenne} d'une classe~:
$\frac{9 \times 10 + 110}{10} = 20$ étudiants. Taille \omterm{def:g11:stat:mean}{moyenne}
vécue par les étudiants~:
$\frac{90 \times 10 + 110 \times 110}{200} = \frac{13\,000}{200} =
65$ étudiants. La brochure imprime $20$~; trois étudiants sur cinq
siègent dans la classe de $110$ et vivent le $65$.

\textbf{16.} Les personnes populaires figurent sur de nombreuses
listes d'amis, si bien que tirer «~un ami~» est biaisé vers les
sociables --- vos amis sont tirés dans les grands \omterm{def:g10:numbers:interval}{intervalles} de
l'horaire social.

\textbf{17.} Pour $t \geq 0$~:
$\P(T > t) = \P(-10\ln U > t) = \P\left(U <
\eu^{-t/10}\right) = \eu^{-t/10}$~: exactement la \omterm{def:g11:func:function}{fonction} de
survie de l'exponentielle --- un seul \omterm{def:g12:exp:ln}{logarithme} convertit le bruit
uniforme de l'ordinateur en n'importe quel temps d'attente.

\textbf{18.} Chaque uniforme a pour \omterm{def:g11:stat:mean}{moyenne} $\frac12$ et pour
\omterm{def:g12:randvar:exp}{variance} $\frac{1}{12}$~: la somme de douze a pour \omterm{def:g11:stat:mean}{moyenne} $6$ et
pour \omterm{def:g12:randvar:exp}{variance} $1$ --- la recette produit donc une \omterm{def:g11:stat:mean}{moyenne} $0$ et
une \omterm{def:g12:randvar:exp}{variance} $1$. Douze petites poussées \omterm{def:g12:sums:indep}{indépendantes}
additionnées~: c'est le mécanisme de Galton, et par la grâce de De
Moivre--Laplace l'histogramme a déjà l'allure d'une cloche à
l'œil nu.

\textbf{19.} Sur le modèle~: les marchés ne sont pas des sommes de
nombreuses petites causes \emph{\omterm{def:g12:sums:indep}{indépendantes}} --- les paniques
corrèlent tout (la leçon de 2008 du problème précédent) et
produisent des queues épaisses qu'aucune \omterm{def:g12:contdist:density}{densité} normale ne
possède. Quand les données murmurent $10^{-137}$, le catéchisme du
statisticien répond~: \emph{c'est le modèle qui est rejeté, pas la
journée}.

\textbf{20.} Uniforme~: toutes les positions se valent, l'honnête
«~je ne connais que les bornes~». Exponentielle~: le risque qui ne
vieillit jamais --- désintégrations, arrivées, déclics. Normale~:
la cloche démocratique des nombreuses petites causes \omterm{def:g12:sums:indep}{indépendantes}
--- tailles, erreurs, \omterm{def:g11:stat:mean}{moyennes}. Piqûre~: ce que vous échantillonnez
est biaisé par la façon dont vous l'avez rencontré --- bus, classes,
amis. Et l'arc~: l'enfant qui comptait des briques de jus de fruits
dans le premier volume a, dix volumes de problèmes plus tard, mesuré
la foule avec une courbe~; l'epsilon, l'\omterm{def:g12:integ:area}{intégrale} et le théorème
central limite attendent dans les volumes universitaires pour
expliquer \emph{pourquoi} la cloche sonne pour tout.
\end{solution}
